\documentclass[11pt,a4paper]{article}
\usepackage[T1]{fontenc}
\usepackage[utf8]{inputenc}
\usepackage{lmodern}
\usepackage[margin=2.5cm]{geometry}
\usepackage{amsmath,amssymb}
\usepackage{booktabs}
\usepackage[hidelinks]{hyperref}
\usepackage{microtype}

\newcommand{\sqT}{\sqrt{T}}
\newcommand{\file}[1]{\texttt{\detokenize{#1}}}

\title{Measured optimality gap of published Clifford+$\sqT$ single-qubit circuits\\[2pt]
\large Technical note}
\author{Cristian Vitorino Rodrigues da Cunha\\
\small Independent researcher, Uberl\^andia, MG, Brazil\\
\small ORCID \href{https://orcid.org/0009-0009-4465-6005}{0009-0009-4465-6005}}
\date{15 November 2026}

\begin{document}
\maketitle

\noindent\textit{Data:} the two witness files are attached to this note. Code, logs and
pre-registration files are listed in the last section and are available on request.

\section*{Summary}

We ran an independent exhaustive search over Clifford+$\sqT$ circuits for the Haar-random targets
published with cyclosynth (Weiden et al., arXiv:2609.16659; data file \file{u3.csv}). Cost is the
authors' own: $\text{cost} = n_T + 3\,n_{\sqT}$.

\begin{itemize}
\item At $\varepsilon = 10^{-3}$, 22 of 500 targets (4.4\%, 95\% Wilson interval 2.9\%--6.6\%) admit a
  strictly cheaper circuit within the same tolerance; the other 478 published circuits are certified
  cost-optimal.
\item At $\varepsilon = 10^{-4}$, 29 of 150 targets (19.3\%, interval 13.8\%--26.4\%) admit a cheaper
  circuit; 120 are certified optimal and 1 is undetermined.
\item The total excess cost is small: 0.29\% at $10^{-3}$ and 1.20\% at $10^{-4}$.
\item All 51 cheaper circuits are listed explicitly and were re-verified through an independent code
  path.
\end{itemize}

The paper states that no cost-optimality is proved, so this note does not contradict it. It measures
the size of the gap and how it grows with precision.

\section*{Method}

Meet-in-the-middle exhaustive search. Every circuit is written as
(canonical word $L$)$\cdot$(canonical word $R$)$\cdot$(Clifford), with $\text{cost}(L) \le a$ and
$\text{cost}(R) \le b$; this covers all circuits of cost $\le a + b - 2$. Canonical words follow
Forest et al.\ (arXiv:1501.04944).

\begin{center}
\begin{tabular}{llllc}
\toprule
Config & $\varepsilon$ & Targets & $a,\ b$ & Covers cost $\le$ \\
\midrule
A & $10^{-3}$ & 0--499 & 12, 11 & 21 \\
B & $10^{-4}$ & 0--149 & 16, 13 & 27 \\
\bottomrule
\end{tabular}
\end{center}

\begin{itemize}
\item Target: $V = R_z(\alpha)\,R_y(\beta)\,R_z(\gamma)$, as in the authors' \file{gather_u3.py}.
\item Acceptance: $d = \sqrt{1 - |\mathrm{tr}(U^\dagger V)|^2/4} < \varepsilon$, strict, no slack.
\item Every circuit found is rebuilt as a gate sequence, its matrix is recomputed exactly in
  $\mathbb{Z}[\zeta_{16}]$, and its cost and distance are recomputed from that matrix.
\item The design was written down before the search code (\file{PREREGISTRO_L3.md}). The search has
  no tunable parameter that changes the outcome.
\end{itemize}

Controls, all passed in both configurations:

\begin{enumerate}
\item The authors' \file{distance} column is reproduced from their gates (relative error
  $\le 10^{-4}$; 0 failures).
\item For every target within reach, the search finds a circuit of cost $\le$ the published cost
  (0 failures).
\item The number of generated words matches the counting recurrence (155{,}596 and 403{,}660).
\item The 24 Cliffords are distinct up to phase and have cost 0.
\end{enumerate}

\section*{Results}

\begin{center}
\begin{tabular}{lll}
\toprule
 & A ($\varepsilon = 10^{-3}$) & B ($\varepsilon = 10^{-4}$) \\
\midrule
Targets & 500 & 150 \\
Cheaper circuit exists & 22 (4.4\%) & 29 (19.3\%) \\
Published circuit certified optimal & 478 & 120 \\
Undetermined & 0 & 1 (target 52) \\
Mean saving where suboptimal ($T$ states) & 1.23 (max 3) & 1.62 (max 4) \\
Excess over total cost & 0.29\% & 1.20\% (47 of 3922) \\
\bottomrule
\end{tabular}
\end{center}

Independent verification (\file{l3_conferir.py}): cost computed with the authors' unmodified
\file{_cost.py}; distance recomputed with 60-digit arithmetic (mpmath), not sharing code with the
search. 22/22 and 29/29 witnesses confirmed; largest $d/\varepsilon$ is 0.998 (A) and 0.975 (B).

Witness files: \file{l3_A_testemunhas.csv}, \file{l3_B_testemunhas.csv}. Columns, with no header
row: $\varepsilon$, target index, published cost, new cost, distance of the cheaper circuit, cheaper
gates, published gates.

\section*{Where the cheaper circuits are (exploratory, not pre-registered)}

In all 51 witnesses the cheaper circuit has a larger denominator exponent than the published one
(mean $+5.7$ levels, range $+1$ to $+14$), fewer $\sqT$ syllables (mean $-2.7$ in A, $-3.0$ in B)
and more $T$ syllables ($+7.0$, $+7.4$). They are deeper and richer in $T$.

\section*{Does a wider search window in cyclosynth recover them?\\
\normalsize (pre-registered, \file{PREREGISTRO_L4.md})}

{\raggedright cyclosynth at commit \texttt{0b1ffd0}, called as the authors do
(\file{sqrt_t=True}, \file{near_clifford_escape=False}). Control: window 2 with no deadline reproduces the
published cost in 149 of 150 targets at $\varepsilon = 10^{-4}$.\par}

On the 29 suboptimal targets at $\varepsilon = 10^{-4}$ (published mean cost 27.79; optimal 26.17),
no deadline:

\begin{center}
\begin{tabular}{ccccc}
\toprule
\file{lde_window} & At the optimum & Improved & Mean cost & Seconds per target \\
\midrule
4 & 8 (28\%) & 9 & 27.17 & 7.9 \\
6 & 9 (31\%) & 11 & 27.07 & 16.3 \\
8 & 12 (41\%) & 16 & 26.86 & 26.0 \\
\bottomrule
\end{tabular}
\end{center}

Window 8 recovers 57\% of the cost gap at about twelve times the run time of window 2 (26.0\,s
against 2.1\,s per target). No returned circuit was cheaper than our optimum, which cross-checks the
exhaustive search.

At $\varepsilon = 10^{-7}$ (20 targets, optimum unknown, published mean 50.20), longer deadlines give
49.20 (window 2, 20\,s), 48.90 (window 4, 20\,s) and 48.70 (window 4, 60\,s): up to 2.99\% below the
published cost, still decreasing with time. With the default settings our own run gave 50.75, which
is 1.10\% above the published cost (8 of 20 targets worse, 4 better), so the published runs are not
reproduced exactly on our machine at this precision. With $n = 20$ this is suggestive only.

Note: the source at this commit already contains a certified lower bound on cost from the denominator
exponent and an internal exhaustive certified mode not exposed in the Python interface. The gap
measured here is therefore a time-versus-cost trade-off in the default settings, not a missing
capability.

\section*{A counting estimate for the slope (exploratory)}

Mooney, Hill and Hollenberg (Quantum 5, 396 (2021), arXiv:2005.05581, Appendix A, Eq.~20) count the
canonical sequences of cost $\le C$ built from $z$-rotations of Clifford-hierarchy levels
$\ell \ge 3$, with $|T_\ell| = 2^{\ell-2}$ rotations of cost $c_\ell$ per level. They use this count
to predict gate proportions. Its asymptotic growth rate is $x^C$, where $x$ is the root of
\[
  \sum_{\ell} 2^{\ell-2}\, x^{-c_\ell} = 1 .
\]
Combined with the standard volume argument (about $\varepsilon^{-3}$ distinct unitaries are needed),
this gives the typical-target estimate
\[
  \text{cost} \approx \frac{3}{\log_2 x}\,\log_2(1/\varepsilon),
\]
with no free parameter. We did not find this step in that paper, where the slopes are reported as
empirical fits; we regard it as a short corollary of their model, not a new result.

Check against their fitted slopes (per $\log_{10}(1/\varepsilon)$, their Tables 1 and 3;
\file{l5_lei_cambio.py}):

\begin{center}
\begin{tabular}{lcc}
\toprule
Gate set & Direct: predicted / measured & Magic states: predicted / measured \\
\midrule
Set 1 & 9.97 / $10.46 \pm 0.43$ & 9.97 / $10.46 \pm 0.43$ \\
Set 2 & 6.62 / $6.89 \pm 0.22$ & 7.25 / $7.47 \pm 0.15$ \\
Set 3 & 5.86 / $6.05 \pm 0.03$ & 6.96 / $7.19 \pm 0.12$ \\
Set 4 & 5.40 / $5.33 \pm 0.06$ & 6.89 / $7.21 \pm 0.25$ \\
Set 5 & 5.02 / $4.84 \pm 0.21$ & 6.88 / $7.21 \pm 0.25$ \\
\bottomrule
\end{tabular}
\end{center}

All ten agree within 5\%. In two cases (Sets 4 and 5, direct) the measured slope is below the
estimate, by 1.3\% and 3.5\%, so the estimate should not be read as a strict bound on a fitted slope
over a finite range.

For Clifford+$\sqT$ with cost 3 per $\sqT$, the same equation is $x^3 = 2x^2 + 4$, $x = 2.594$,
giving $2.18\,\log_2(1/\varepsilon)$, against the $2.4\,\log_2(1/\varepsilon)$ fitted in
arXiv:2609.16659. Our optimal costs at $10^{-3}$ and $10^{-4}$ (18.71 and 25.83 on different target
sets) correspond to about 2.14 per bit between the two precisions.

\section*{Limits}

\begin{itemize}
\item Exhaustive search does not reach $\varepsilon \le 10^{-5}$; the growth of the gap beyond
  $10^{-4}$ is not tested.
\item Configuration B covers 150 targets; one is undetermined.
\item The counting section and the mechanism section are exploratory. The percentage reductions of
  Mooney et al.\ had been seen before the estimate was computed; the costs and slopes had not.
\item Literature check for the slope estimate was limited to arXiv:2005.05581, 1501.04944,
  2203.10064, 2609.16659 and two web searches (15 Nov 2026). It may exist elsewhere.
\item Run times were measured on a 12-thread laptop under WSL and are not comparable to the paper's.
\end{itemize}

\section*{Files}

Attached: \file{l3_A_testemunhas.csv}, \file{l3_B_testemunhas.csv}.

{\raggedright\noindent Available on request: \file{PREREGISTRO_L3.md}, \file{l3.py}, \file{l3_conferir.py},
\file{l3_A.log}, \file{l3_B.log}, \file{PREREGISTRO_L4.md}, \file{l4_rodar.py},
\file{l4_analise.py}, \file{l4_A.csv}, \file{l4_B.csv}, \file{l5_lei_cambio.py},
\file{l5_lei_cambio.log}.\par}

\end{document}
